%! AP Statistics — Unit 1, Lesson 1.7 — Guided Notes ANSWER KEY
\documentclass[10pt]{article}
\usepackage{apstats-article}
\usepackage{apstats-key}

%=============================================================================
\begin{document}

\pageheader{Unit 1, Lesson 1.7}{Guided Notes --- Answer Key}
\begin{tabularx}{\linewidth}{X X X}
Name:\;\ans{Key} & Date:\; & Period:\;
\end{tabularx}

\vspace{0.10in}

% ── OBJECTIVE ────────────────────────────────────────────────────────────────
\begin{objectivebox}
\small By the end of this lesson, I will be able to\dots\\[4pt]
\ansline{Calculate mean, median, range, IQR, and standard deviation for a
quantitative dataset and interpret each statistic in context.}\\
\ansline{Apply the $1.5\times$IQR fence rule to identify outliers.}\\
\ansline{Choose median + IQR vs.\ mean + $s_x$ based on the shape of the distribution.}
\end{objectivebox}

\vspace{0.10in}

% ── VOCABULARY ───────────────────────────────────────────────────────────────
\begin{vocabbox}
\small
\noindent\textbf{\textcolor{navy}{Mean ($\bar{x}$):}}\\
\ans{The arithmetic average: sum of all values divided by $n$. Not resistant to outliers.}\\[6pt]

\noindent\textbf{\textcolor{navy}{Median ($M$):}}\\
\ans{The middle value of a sorted dataset; average of the two middle values when $n$ is even. Resistant to outliers.}\\[6pt]

\noindent\textbf{\textcolor{navy}{Range:}}\\
\ans{max $-$ min; simple but not resistant since it uses only the two extreme values.}\\[6pt]

\noindent\textbf{\textcolor{navy}{Quartiles ($Q_1$, $Q_3$):}}\\
\ans{$Q_1$ = median of the lower half of the data (25th percentile); $Q_3$ = median of the upper half (75th percentile).}\\[6pt]

\noindent\textbf{\textcolor{navy}{IQR (Interquartile Range):}}\\
\ans{$Q_3 - Q_1$; the spread of the middle 50\% of data. Resistant to outliers. Pair with median.}\\[6pt]

\noindent\textbf{\textcolor{navy}{Resistant statistic:}}\\
\ans{A statistic not strongly affected by extreme values; the median and IQR are resistant, the mean and standard deviation are not.}\\[6pt]

\noindent\textbf{\textcolor{navy}{Outlier fence rule:}}\\
\ans{A value is an outlier if it falls below $Q_1 - 1.5\cdot\text{IQR}$ (lower fence) or above $Q_3 + 1.5\cdot\text{IQR}$ (upper fence).}
\end{vocabbox}

\vspace{0.10in}

% ── HOOK ─────────────────────────────────────────────────────────────────────
\begin{hookbox}
\small\textbf{Which salary best represents ``typical''?}

\vspace{5pt}
Salaries (thousands): \$30, \$32, \$35, \$38, \$200

\vspace{4pt}
\renewcommand{\arraystretch}{1.4}
\begin{tabularx}{\linewidth}{@{} >{\bfseries}p{0.20\linewidth} X X @{}}
\rowcolor{sky}
\textbf{Statistic} & \textbf{Calculation} & \textbf{Interpretation} \\
\midrule
Mean   & \ans{$335/5 = \$67$k} & \ans{pulled up by the \$200k outlier} \\[8pt]
Median & \ans{\$35k (middle value)} & \ans{resistant; typical worker's salary} \\
\end{tabularx}

\vspace{6pt}
Which is more representative? \ans{Median (\$35k).}
Why? \ans{Mean is distorted by the extreme value of \$200k.}
\end{hookbox}

\vspace{0.08in}

% ── SECTION 1: MEASURES OF CENTER ────────────────────────────────────────────
\begin{notesbox}{Measures of Center}
\small
\begin{multicols}{2}

\textbf{\textcolor{navy}{Mean} ($\bar{x}$):}
\[\bar{x} = \frac{\sum x_i}{n}\]
\begin{itemize}[itemsep=3pt]
  \item Uses \ans{all} data values.
  \item \ans{Not} resistant to outliers.
  \item Pulled toward the \ans{tail} of a skewed distribution.
\end{itemize}

\columnbreak

\textbf{\textcolor{navy}{Median} ($M$):}
\begin{itemize}[itemsep=3pt]
  \item Sort data; take the \ans{middle} value.
  \item If $n$ is even, average the \ans{two} middle values.
  \item \ans{Is} resistant to outliers.
  \item Stays near the \ans{center} of the distribution.
\end{itemize}
\end{multicols}
\end{notesbox}

\vspace{0.08in}

% ── SECTION 2: MEASURES OF SPREAD ────────────────────────────────────────────
\begin{notesbox}{Measures of Spread}
\small
\begin{multicols}{2}

\textbf{\textcolor{navy}{Range:}}
\[\text{Range} = \max - \min\]
\begin{itemize}[itemsep=3pt]
  \item Simple but \ans{not} resistant.
  \item Only uses the two \ans{extreme} values.
\end{itemize}

\vspace{6pt}
\textbf{\textcolor{navy}{Quartiles \& IQR:}}
\begin{itemize}[itemsep=3pt]
  \item $Q_1 = $ median of the \ans{lower} half.
  \item $Q_3 = $ median of the \ans{upper} half.
  \item $\text{IQR} = Q_3 - Q_1$
  \item Contains the middle \ans{50}\% of data.
  \item \ans{Is} resistant to outliers.
\end{itemize}

\columnbreak

\textbf{\textcolor{navy}{Standard Deviation} ($s_x$):}
\[s_x = \sqrt{\frac{\sum(x_i - \bar{x})^2}{n-1}}\]
\begin{itemize}[itemsep=3pt]
  \item Typical \ans{distance} from the mean.
  \item Always $\geq$ \ans{0}; units = units of data.
  \item \ans{Not} resistant to outliers.
  \item Use with \ans{mean} when distribution is symmetric.
\end{itemize}
\end{multicols}
\end{notesbox}

\vspace{0.08in}

% ── SECTION 3: OUTLIER RULE + DECISION ───────────────────────────────────────
\begin{notesbox}{Outlier Rule \& Choosing Statistics}
\small
\begin{multicols}{2}

\textbf{\textcolor{navy}{$1.5 \times$ IQR Fence Rule:}}
\begin{itemize}[itemsep=4pt]
  \item Lower fence: $Q_1 - 1.5(\text{IQR})$
  \item Upper fence: $Q_3 + 1.5(\text{IQR})$
  \item A value is an outlier if it is \ans{below} \\
        the lower fence or \ans{above} the upper fence.
\end{itemize}

\vspace{4pt}
\textbf{Example:} $Q_1 = 19$, $Q_3 = 31$, IQR~$= 12$\\
Lower fence $= \ans{1}$\quad Upper fence $= \ans{49}$

\columnbreak

\textbf{\textcolor{navy}{Choosing the Right Pair:}}

\vspace{4pt}
\renewcommand{\arraystretch}{1.5}
\begin{tabularx}{\columnwidth}{@{} >{\bfseries}p{0.42\columnwidth} X @{}}
\rowcolor{sky}
\textbf{If shape is\dots} & \textbf{Report\dots} \\
\midrule
Skewed \textit{or} has outliers & \ans{Median + IQR} \\
Approximately symmetric         & \ans{Mean + $s_x$} \\
\end{tabularx}

\vspace{6pt}
\textit{Key:} Mean is pulled toward the \ans{tail}. \\
Median \ans{resists} toward extremes.
\end{multicols}
\end{notesbox}

\vspace{0.08in}

% ── REFLECTION ───────────────────────────────────────────────────────────────
\begin{tcolorbox}[colback=goldbg, colframe=goldacc, boxrule=0.8pt, arc=2mm,
    left=4mm, right=4mm, top=2mm, bottom=2mm]
\small\textbf{Quick Check:} A dataset has mean 85 and median 72.\\[4pt]
Shape is likely \ans{skewed right}. \quad
Better center: \ans{median}. \quad
Better spread: \ans{IQR}.
\end{tcolorbox}

\end{document}
